\section{Amplifying an above-half-rate seed}\label{sec:amplification}

The construction uses a fixed finite field of odd characteristic; its
finite operator only needs $2\ne0$ in the field. All
logarithms in the decay constants of this section are natural logarithms;
binary rates and label budgets use $\log_2$.

\begin{definition}[The seed hypothesis]\label{def:seed}
\formal{TapeStorageRates/OddGirthLinearSeed.lean}{TapeStorageRates.OddGirthLinearSeed}
An \emph{odd-girth linear seed} over a finite field $\mathbb F_q$, with
degree bound $D$ and rate parameter $\rho$, consists of simple graphs
$G_j$ on $n_j$ coordinates and linear subspaces
$U_j\le\mathbb F_q^{n_j}$ such that $n_j\to\infty$, every degree is at
most $D$, $U_j$ is a storage code for the neighbourhood relation of
$G_j$, and $\dim U_j\ge\rho n_j$. For every fixed $r$, eventually every
odd closed walk in $G_j$ has length greater than $r$.
\end{definition}

\begin{proposition}[A concrete seed]\label{prop:seed-exists}
\notformalized
There is an odd-girth linear seed with
$q=127^2$, $\rho=810/1613>1/2$, and a fixed finite degree bound $D$.
\end{proposition}
\begin{informalproof}
Use the graphs and linear codes of \cite[Corollary~U(a,b)]{PaperThirteen}.
Write $\Delta_{13}$ for the fixed incidence-degree bound of that paper's
Theorem~A, and set $D=2\Delta_{13}<2\cdot127^{8\,791\,448}$.
Discard the finite initial set of
levels for which the displacement bound does not yet exceed the check
radius. Corollary~U(a) gives simple graphs of maximum degree at most
$D$. Every odd closed walk at level $j$ has length at least
\[
 \frac{j\log509-\log4}{a\log127},\qquad a=2{,}197{,}862,
\]
which tends to infinity. Part~(b) supplies an $\mathbb F_{127^2}$-linear
storage code on the same graph with dimension at least
$(8100/16130)n_j=(810/1613)n_j$. Theorem~A's construction has
$n_j=16130c_j\to\infty$. Relabel each finite coordinate set by
$\{0,\ldots,n_j-1\}$ and reindex the remaining levels. Relabelling
preserves dimension, recovery, degrees and walk lengths, giving every
condition of \cref{def:seed}. The field has odd characteristic, and
$1/2<810/1613<1$.

This application uses the written arithmetic-geometric existence proof
in the cited paper. It does not assert that the concrete seed has been
constructed in Lean. The results below take precisely
\cref{def:seed} as their existence hypothesis.
\end{informalproof}

On a linear code, deterministic recovery from a set of coordinates
can be represented by a linear functional on those coordinates.
Consequently each seed level has a matrix $A$ supported on the edges
of $G_j$, with $Ax=x$ for $x\in U_j$. This adapter is
\lean{TapeStorageRates.OddGirthLinearSeed.exists_recovery_matrix}, reusing
\lean{LeanMachine.LinearRecovery.IsRecoverable.exists_recovery_matrix}.
No linear decoder is an additional seed assumption.

\paragraph{The finite operator.}
For a matrix $A$ on $V$ over a field $K$, an odd integer $m\ge1$, and
$\sigma\in\{-1,1\}^m$, let $\operatorname{Maj}(\sigma)$ be the sign
occurring more often. Take this majority before interpreting the signs
in the field. Define, on $(K^V)^{\otimes m}$,
\begin{equation}\label{eq:majority}
 B_m=2^{-m}\sum_{\sigma\in\{-1,1\}^m}
       \operatorname{Maj}(\sigma)
       \bigotimes_{i=1}^m(I+\sigma_i A),
 \qquad H_m=I-B_m.
\end{equation}
Matrices on $V^m$ always use the standard coordinate tensor basis.
For $1/2<\rho<1$, put
\begin{equation}\label{eq:decay}
 \lambda=-\log\bigl(2\sqrt{\rho(1-\rho)}\bigr)
         =-\tfrac12\log\bigl(4\rho(1-\rho)\bigr)>0.
\end{equation}

\begin{theorem}[Theorem M: finite majority amplification]\label{thm:majority}
\formal{TapeStorageRates/MajorityRank.lean}{TapeStorageRates.Majority.check_rank_add_binomial_le,TapeStorageRates.Majority.check_rank_le_exp}
Let $G$ be a simple graph on $N>0$ vertices, with maximum degree at
most $D$. Let $A$ be a matrix over a field of characteristic different
from two, supported on edges of $G$, and let $U\le\ker(I-A)$ have
dimension $h$. For every odd $m\ge1$, the matrix $H_m$ in
\cref{eq:majority} has diagonal one and satisfies
\begin{equation}\label{eq:exact-rank}
 \operatorname{rank}H_m+
 \sum_{\substack{0\le a\le m\\2a>m}}\binom ma h^a(N-h)^{m-a}\le N^m.
\end{equation}
If $h\ge\rho N$ with $1/2<\rho<1$, then
$\operatorname{rank}H_m\le e^{-\lambda m}N^m$.
Define $N_{H_m}(v)=\{w\ne v:(H_m)_{vw}\ne0\}$.
Every port in its incidence graph has degree at most $(D+1)^m-1$.
If every odd closed walk in $G$ has length greater than $r$, then
every paired distance for $N_{H_m}$ is greater than $r$.
\end{theorem}
\begin{outline}
Choose a basis of $U$ and extend it by an arbitrary complementary basis.
A pure tensor using more than $m/2$ basis vectors from $U$ is fixed by
$B_m$: any corner with negative majority has a negative sign at a fixed
factor and therefore kills the tensor. The sum of all unsigned corners
is $2^mI$. The independent tensor basis vectors with a strict majority
of fixed factors number exactly the sum in \cref{eq:exact-rank}.
Rank--nullity gives the inequality. The complement need not be
$A$-invariant. This is checked by
\lean{TapeStorageRates.Majority.operator_tprod_eq_self_of_majority} and
\lean{TapeStorageRates.Majority.subspace_goodIndices_card_le_finrank_ker}.

For the exponential estimate, put $t=h/N\ge\rho$. The fraction of
excluded basis vectors is the binomial lower tail
$\sum_{a=0}^{\lfloor m/2\rfloor}\binom ma t^a(1-t)^{m-a}$.
For $t\ge1/2$ and $2a\le m$, each product is at most
$[t(1-t)]^{m/2}$. Summing the binomial coefficients gives
$[2\sqrt{t(1-t)}]^m\le e^{-\lambda m}$.
\lean{LeanMachine.BinomialTail.half_sum_le_exp} proves this finite-sum
bound, including the endpoint $t=1$.

For the support statement, $A$ has zero diagonal. Negating all signs
in \cref{eq:majority} negates their majority. Pairing these sign patterns
shows that a nonzero entry of $B_m$ changes an odd number of coordinates;
each changed coordinate follows an edge of $G$. The stationary tuple
is excluded, so the containing graph has degree at most $(D+1)^m-1$.
An odd closed walk in this graph makes an odd total number of coordinate
changes. Some coordinate projection, after stationary steps are removed,
is an odd closed walk in $G$ of no greater length. Finally a paired
incidence walk projects to an odd closed walk of this kind. The support,
degree and distance conclusions are
\lean{TapeStorageRates.Majority.check_diagonal},
\lean{TapeStorageRates.Majority.check_helpers_degree_add_one_le} and
\lean{TapeStorageRates.Majority.check_helpers_paired_far}.
\end{outline}

\paragraph{Binary conversion on the same coordinates.}
For a matrix $H$ over $\mathbb F_q$, embed each input bit as $0$ or $1$
and let $\sigma_H(x)=Hx$. The nominal syndrome has $N$ field entries,
but its attained labels lie in an image of dimension $\operatorname{rank}H$.
We charge that image size, not the nominal output width.

\begin{proposition}[Syndrome fibres]\label{prop:syndrome}
\formal{TapeStorageRates/SyndromeRates.lean}{TapeStorageRates.matrixSyndrome_subset_fibre_storage,TapeStorageRates.matrixSyndrome_image_card_le_rpow,TapeStorageRates.matrix_exists_binary_code_rate,TapeStorageRates.matrixSyndrome_binary_rate_tail}
Let $H$ be an $N\times N$ matrix over $\mathbb F_q$, with $N>0$ and
nonzero diagonal, and put $h=\operatorname{rank}H$.
The map $\sigma_H:\{0,1\}^N\to\mathbb F_q^N$ has at most $q^h$
attained labels. Every nonempty intersection of a fibre with a source
subset is a binary storage code for
$N_H(v)=\{w\ne v:H_{vw}\ne0\}$.
Some fibre has rate at least $1-h\log_2q/N$.
For every nonempty $S\subseteq\{0,1\}^N$, write
$S_x=\{y\in S:\sigma_H(y)=\sigma_H(x)\}$ and
$r(S)=\log_2|S|/N$. For every real $u$,
\begin{equation}\label{eq:syndrome-tail}
 \Pr_{x\ \mathrm{uniform\ on}\ S}
 \left[r(S_x)<r(S)-\frac{h\log_2q}{N}-u\right]
 \le 2^{-uN}.
\end{equation}
The same map and helper relation work for every $S$ and $u$.
\end{proposition}
\begin{outline}
If two words have the same syndrome and agree on $N_H(v)$, subtracting
their $v$th equations gives $H_{vv}(x_v-y_v)=0$. The nonzero diagonal
and the injective binary embedding give $x_v=y_v$.
The linear image contains $q^h$ elements, so pigeonhole counting gives
a fibre of size at least $2^N/q^h$. More generally, partition $S$ by its
attained labels. For each attained cell, summing $1/|S_x|$ over its
members gives one. Hence
\[
 \sum_{x\in S}2^{\log_2|S|-\log_2|S_x|}
     =|S|\,|\sigma_H(S)|\le |S|q^h.
\]
Counting the terms above $q^h2^{uN}$ gives
\cref{eq:syndrome-tail}. The exact identity is
\lean{TapeStorageRates.matrixSyndrome_deficit_moment}; it is an instance
of the same finite-partition identity used for the native tail theorem.
\end{outline}

\begin{theorem}[Theorem C: every sufficiently large degree cap]\label{thm:construction}
\formal{TapeStorageRates/MajorityPartitions.lean}{TapeStorageRates.OddGirthLinearSeed.exists_majority_partition_family,TapeStorageRates.OddGirthLinearSeed.exists_partition_family_at_cap}
Assume an odd-girth linear seed over $\mathbb F_q$ of odd characteristic,
with $1/2<\rho<1$. With $\lambda$ as in \cref{eq:decay}, set
\begin{equation}\label{eq:construction-constants}
 C=(\log_2q)e^{2\lambda},\qquad
 \alpha=\frac{\lambda}{\log(D+1)},\qquad
 \varepsilon_\Delta=C\Delta^{-\alpha}.
\end{equation}
Here $C,\alpha>0$. For each integer $\Delta\ge D+1$, there are positive
lengths $n_j\to\infty$, helper relations $N_j$, and maps
$f_j:\{0,1\}^{n_j}\to Y_j$ such that every incidence degree is at
most $\Delta$, the minimum paired distance tends to infinity, and
\[
 |f_j(\{0,1\}^{n_j})|\le2^{\varepsilon_\Delta n_j}
 \qquad\text{for every }j.
\]
Each attained fibre is a storage code for $N_j$. In particular, there
is an admissible binary family at cap $\Delta$ whose rate is at least
$1-\varepsilon_\Delta$ at every index.
\end{theorem}
\begin{outline}
A positive-dimensional seed forces $D>0$. Fix an odd amplification
order $m$ before the seed level grows. Apply \cref{thm:majority} to
the supported recovery matrices at every level, and then
\cref{prop:syndrome}. The resulting partitions have label exponent
$(\log_2q)e^{-\lambda m}$. Their sizes grow, and every fixed paired
distance threshold is inherited from the seed. A finite initial prefix
is discarded to make every size positive.

For $\Delta\ge D+1$, choose the largest odd $m$ with
$(D+1)^m\le\Delta$. Then
$m>\log\Delta/\log(D+1)-2$, which gives
$e^{-\lambda m}\le e^{2\lambda}\Delta^{-\alpha}$.
The discrete choice and estimate reuse
\lean{LeanMachine.OddPowerChoice.exists_odd_power_and_decay}.
The field, seed and constants precede the degree cap; $m$ precedes
the family index. Choosing one largest fibre at each level gives
\lean{TapeStorageRates.OddGirthLinearSeed.exists_family_at_cap}.
\end{outline}

The theorem is an existence construction. Its recovery matrices and
syndrome labels have not been supplied with an efficient uniform
generation or evaluation algorithm. Its relevance is that the abstract
rate requirement of Theorem~A already holds without a transposer.
